% TOP_PROBLEM: {"id": 3341, "title": "MSO alternation hierarchy over pictures", "category_id": 18, "category": {"id": 18, "name": "logic", "display_name": "Logic", "description": "Problems in mathematical logic, model theory, proof theory, and finite model theory.", "slug": "logic", "order_index": 18, "created_at": "2026-07-31T15:26:25.671Z"}, "status": "open", "problem_number": "OPG-37448", "set_id": 12}
% FIRST_POSTED: 2026-10-02
% ATTEMPT_STATUS: unresolved
% MODEL: GPT 6 Astra Ultra
% UnsolvedMath ID 3341; source catalogue code OPG-37448
% !TeX program = lualatex
% Research attempt dated 2026-10-02; canonical statement preserved.
\documentclass[11pt,a4paper]{article}
\usepackage[margin=25mm]{geometry}
\usepackage{fontspec}
\setmainfont{lmroman10-regular.otf}[BoldFont=lmroman10-bold.otf,
 ItalicFont=lmroman10-italic.otf,BoldItalicFont=lmroman10-bolditalic.otf]
\usepackage{amsmath,amssymb,mathtools}
\usepackage{unicode-math}
\setmathfont{latinmodern-math.otf}
\AtBeginDocument{\let\setminus\smallsetminus}
\usepackage{xurl,hyperref}
\hypersetup{colorlinks=true,urlcolor=blue}
\setcounter{secnumdepth}{0}
\setlength{\parindent}{0pt}
\setlength{\parskip}{5pt}
\setlength{\emergencystretch}{3em}
\begin{document}
\section*{MSO alternation hierarchy over pictures}\label{3341}
{\small UnsolvedMath ID 3341 · OPG-37448 · Logic · OpenGarden}
\subsection{Definitions and mathematical statement}
A picture over a fixed finite alphabet $A$ is a map
$P:[h]\times[w]\to A$, where $h,w\ge1$ and $[h]=\{1,\ldots,h\}$.
Regard it as a finite relational structure with horizontal and vertical
successor relations and one unary predicate for each letter. Monadic
second-order logic (MSO) permits quantification over cells and over sets of
cells. A $\Sigma_k$ MSO sentence has at most $k$ alternating blocks of
set quantifiers, beginning existentially, followed by a first-order formula;
first-order quantifiers inside that formula do not count as set alternations.

Fix a class $\mathcal B$ of balanced pictures. Strictness asks, for every
$k\ge1$, for a language $L\subseteq\mathcal B$ definable by a
$\Sigma_{k+1}$ sentence but by no $\Sigma_k$ sentence on $\mathcal B$.
The source asks this for linearly and polynomially balanced rectangles.
To make the balance uniform, take fixed constants $C\ge1,d\ge1$ and require
\[
 \max(h,w)\le C\min(h,w)^d.
\]
Investigate strictness for these fixed classes, separating $d=1$ from
$d>1$; square pictures $h=w$ are the basic special case. This explicit
family of balance conventions resolves the source's unspecified polynomial.
Allowing a different bound for each picture would impose no restriction.
Equivalence of sentences is tested only on the chosen class of pictures.
\subsection{Short English statement}
Does each additional alternation of set quantifiers express genuinely new properties even for pictures whose two side lengths remain comparable?
\subsection{Research attempt}
\paragraph{Scope and method.}
This research attempt by GPT-6 Astra Ultra is dated 2 October 2026.
It preserves the catalogue statement and distinguishes elementary proofs
below from cited prior work. No novelty claim or formal proof-assistant
verification is made. The final paragraph states the precise remaining scope.
The finite checks described below are reproducible in
\texttt{scripts/check\_attempt\_batch03\_root\_2026\_10\_02.py}.
\paragraph{A restriction that destroys familiar witnesses.}
The constants $C,d$ are fixed for the class, rather than
chosen separately for each rectangle. The source discussion
[S2] describes known hierarchy witnesses on very unbalanced
rectangles, whose long side grows as an exponential tower
in the short side. Restricting such a witness language to
the balanced class is not a proof of separation on that
class. We make the obstruction precise and isolate a
valid first-level transfer argument.

\paragraph{Finite picture languages are first-order.}
For any particular finite picture with $N$ cells, quantify
first-order variables $x_1,\ldots,x_N$. Require them to be
pairwise distinct and require every cell to equal one of
them. For every ordered pair of these variables, specify
whether the horizontal and vertical successor relations
hold, and specify each cell's letter predicates. This
sentence characterizes the picture up to isomorphism.
A finite disjunction characterizes any finite language
of pictures. It uses no set quantifiers at all. In
particular, every property restricted to a bounded total
number of cells is first-order definable; inspecting a
finite collection of pictures cannot prove a strict
set-quantifier hierarchy.

\paragraph{Applying the lemma to tower-shaped supports.}
Suppose a proposed separator is supported on rectangles
of dimensions $n$ by $f(n)$, where $f(n)\ge2^n$ for
all sufficiently large $n$. Balance would require
$f(n)\le Cn^d$ once $f(n)\ge n$. But
$2^n/(Cn^d)$ tends to infinity for fixed $C,d$.
There are therefore only finitely many admissible values
of $n$. Over a fixed finite alphabet, those finitely many
shapes admit only finitely many pictures. The restricted
language is first-order by the preceding lemma, no matter
how high its unrestricted alternation complexity was.
For the concrete support $f(n)=2^n$, $C=1,d=2$,
the admissible positive integers are exactly $2,3,4$.
At $n=5$ the inequality fails, and the ratio $2^n/n^2$
is strictly increasing for $n\ge3$, since
$2n^2>(n+1)^2$ there. This small example exhibits the
loss of the unbounded witness parameter explicitly.

\paragraph{What complement nonclosure really gives.}
Let $\mathcal B$ be a fixed class of pictures. Assume a
language $L\subseteq\mathcal B$ has an existential MSO
definition $\exists X_1\cdots\exists X_s\,\varphi$, with
$\varphi$ first-order, and that its complement inside
$\mathcal B$ has no such definition. That complement
is defined by $\forall X_1\cdots\forall X_s\,\neg\varphi$.
Prefixing an unused existential set variable puts it in
$\Sigma_2$. Hence this nonclosure hypothesis implies
$\Sigma_1\ne\Sigma_2$ on the class. The source reports
the relevant nonclosure result for squares; this paragraph
proves the logical implication, not that imported result.
The identical transfer works at any level if nonclosure
at that level is known. Nonclosure at level one alone
does not give the required hypothesis at every later level.

\paragraph{The padding gap.}
One might place a thin witness rectangle in a much larger
square and mark the active region. To transfer a lower
bound, however, one must pull every hypothetical
low-alternation definition on the padded objects back
to one on the original rectangles without increasing
the number of set-quantifier blocks. Merely encoding
the rectangle injectively does not supply that translation.
Sets quantified over the large padding can contain many
more independent bits than sets on the original domain.
No suitable alternation-preserving interpretation is
constructed here.

\paragraph{Verification and final status.}
The root script checks the explicit exponential-versus-square
example; the monotonicity and finite-language arguments
prove its tail and its logical consequence. The attempt
establishes why direct restriction fails and gives a
conditional separation lemma. It does not construct a
balanced separator at each level, so the full hierarchy
question remains unresolved here.

\subsection{Sources}
\begin{itemize}
\item[S1] ULAM.AI and the UnsolvedMath Contributors, supplied problems.json, record 3341 (OPG-37448), OpenGarden collection. Catalogue identity and source transcription; CC BY 4.0.
\url{https://huggingface.co/datasets/ulamai/UnsolvedMath}.
\item[S2] Open Problem Garden, MSO alternation hierarchy over pictures. Original problem and discussion; the mathematical formulation below is expanded from the supplied source record.
\url{http://www.openproblemgarden.org/op/mso_alternation_hierarchy_over_pictures}.

\end{itemize}
\end{document}
